% Direct all-dimensional local application; no higher-dimensional resolution is used.
\begin{corollary}[Proved: bounded scalar obstructions on complete local punctures]
\label{D:local-bounded-scalar-simple-product}
Let $(R,\mathfrak m,k)$ be a complete excellent normal local
domain of dimension at least two, equicharacteristic $p>0$, with
algebraically closed residue field. Put
\[
          U=\operatorname{Spec}R\setminus\{\mathfrak m\},
          \qquad G=\pi_1(U).
\]
Fix $B\ge1$. For every $q=p^f$ and every finite subgroup
$\widetilde S\subset\operatorname{GL}_r(\mathbf F_q)$ whose
scalar kernel has order at most $B$, let $S$ be its image in
$\operatorname{PGL}_r(\mathbf F_q)$. There is a constant
$C_{R,B}$ such that
\[
 \#\{\operatorname{Hom}_{\mathrm{cont}}(G,S)
            \text{ up to ambient projective conjugacy}\}
                  \le q^{C_{R,B}r^2}.                  \tag{D.LSC.1}
\]
If a family of nonabelian simple groups admits such realizations
with $|S|\ge q^{\epsilon r^2}$ for one fixed $\epsilon>0$, the
quotient of $G$ formed by all distinct simple quotient kernels
of these types is topologically finitely generated.

In particular this holds for all the nonabelian simple groups
$\operatorname{PSL}_r(p^f)$ with $\gcd(r,p^f-1)\le B$, and
for all $\operatorname{PSU}_r(p^f)$ with
$\gcd(r,p^f+1)\le B$, with unbounded $r$ in either assertion.
The unitary representation count is understood in
$\operatorname{PGL}_r(\mathbf F_{p^{2f}})$.
\end{corollary}
\begin{proof}
Completeness makes $R$ henselian, and the residue field is
algebraically closed, so $R$ is strictly henselian.
Corollary~\ref{D:strict-local-open-cohomology-finite}, which
uses precisely the constructibility theorem
\citep[XIII, Theorem~1.1.1(i)]{A_ILO2014}, proves that
$H^2(G,\mathbf Z/d\mathbf Z)$ is finite for every $d$ prime
to $p$. We use only finiteness for each $d$, not a bound
uniform in $d$.

We verify the hereditary representation hypothesis carefully.
Let $H\subset G$ be open normal, and let $V\to U$ be the
corresponding connected finite etale cover. The normal connected
scheme $U$ is integral, as is $V$. Normalize $R$ in $k(V)$,
and write the resulting finite normal domain as $R_H$.
Finiteness follows from excellence. Since a finite algebra
over a complete henselian local ring is a product of complete
local rings and $R_H$ is a domain, $R_H$ is local and complete.
Its residue field is a finite extension of $k$, hence is $k$.
It is excellent, normal, equicharacteristic $p$, and has the
same dimension as $R$. Its restriction above $U$ is exactly
$V$, by uniqueness of finite normalization, while its closed
fibre has one underlying point. Consequently
\[
       V=\operatorname{Spec}R_H\setminus\{\mathfrak m_H\},
       \qquad H=\pi_1(V).
\]
Corollary~\ref{D:uniform-local-growth-all-dimensions} therefore
supplies a constant $c_H$ with
\[
 \#\bigl(\operatorname{Hom}_{\mathrm{cont}}
       (H,\operatorname{GL}_r(\mathbf F_{p^f}))/\mathrm{conj}\bigr)
                     \le p^{c_Hfr^2}.
\]
This all-dimensional quadratic theorem uses the finite surface
slice and proper normalization descent of
Proposition~\ref{D:local-surface-reduction}; it requires no
resolution in the dimension of $R$.

Thus both hypotheses of
Proposition~\ref{D:bounded-scalar-projective-growth} hold for
$G$, proving (D.LSC.1) and the simple-product assertion.
The indicated linear and unitary families follow from
Corollaries~\ref{D:surface-bounded-centre-simple-product} and
\ref{D:bounded-centre-unitary-simple-product}'s elementary
matrix-cover and group-order calculations; their geometric
surface hypotheses are not being invoked here.
\end{proof}

\begin{remark}[Local scope]
The prime-to-characteristic cohomology input itself holds on
arbitrary-dimensional strictly local quasi-excellent opens.
The quantitative conclusion here uses the separate quadratic
representation theorem for complete equicharacteristic normal
punctures with algebraically closed residue field. We do not
extend that quantitative theorem to mixed characteristic or
to an arbitrary open subscheme of a local scheme.
\end{remark}
